% standardmäßig im 16:9 Format
% verwende \documentclass[43]{sikslides} für 4:3
\documentclass{include/sikslides}
\usepackage{graphicx}
\usepackage{listings}
\usepackage{xcolor}

\lstset{
	basicstyle=\ttfamily\footnotesize,
	keywordstyle=\color{blue},
	commentstyle=\color{gray},
	stringstyle=\color{red},
	showstringspaces=false,
	frame=single,
	breaklines=true
}


\title[Forschungsmodul]{Generierung von Task Graphs}
\subtitle{Forschungsmodul Embedded Systems}
\author{Anja Graf}
\date[01.08.2025]{01. August 2025}

\begin{document}
\titleframe

% erneutes Bauen notwendig, um Gliederung zu aktualisieren
\begin{frame}
    \frametitle{Gliederung}
    \tableofcontents
\end{frame}

\section{Zielsetzung des Forschungsmoduls}

\begin{frame} {\secname }
	\begin{columns}[c]
	\begin{column}{0.6\textwidth}
		\begin{itemize}
			\item Generieren eines Taskgraphen mit randomisierter Graphstruktur
			\begin{itemize}
				\item Dabei: Knoten sind Tasks, Kanten sind Abhängigkeiten 
				\item Graph ist Directed und Acyclic (DAG)
			\end{itemize}
			\item Command Line Interface für Benutzerschnittstelle
			\begin{itemize}
				\item Übergabeparamter wie Knoten- oder Kantenanzahl
			\end{itemize}
			\item Output in Form von generierten Dateien 
		\end{itemize}
	\end{column}
	\begin{column}{0.3\textwidth}
		\includegraphics[width=0.4\linewidth]{images/undirected/graphcircle.png}
		\vspace*{0.1em}
		\centering
		\parbox{\linewidth}{ % Breite der Box (hier halbe Textbreite)
			\raggedleft % Text rechtsbündig innerhalb der Box
			\scriptsize % kleine Schrift
			Ungerichteter, zyklischer Graph \\
			Mit 3 Knoten, 3 Kanten
		}
		\vspace*{0.1em}
		\includegraphics[width=0.6\linewidth]{images/directed/graph.png}
		\vspace*{0.1em}
		\centering
		\parbox{\linewidth}{ % Breite der Box (hier halbe Textbreite)
			\raggedleft % Text rechtsbündig innerhalb der Box
			\scriptsize % kleine Schrift
			Gerichteter, azyklischer Graph \\
			Mit 3 Knoten, 2 Kanten
	}
	\end{column}
	\end{columns}
\end{frame}
\section{Funktionsweise und Inhalt}
\subsection{Erdős–Rényi model}
\begin{frame} {\secname : \subsecname}
	\begin{center}
		\Large\textbf{Zwei Varianten des Modells:}
	\end{center}
	\begin{columns}
		\column{0.55\textwidth} 
		$G(n,p)$ \\
		n = Anzahl an Knoten \\
		p = Wahrscheinlichkeit für Kanten
		\begin{itemize}
			\item Iteriere über geshuffelte Knotenvektoren, um Knoten i, j für jede mögliche Kante $e=\{i,j\} $ zu erhalten ($i \ne j$)
			\item Füge diese Kante e mit Wahrscheinlichkeit p hinzu
		\end{itemize}
	
		
		\column{0.45\textwidth}
		$G(n,m)$ \\
		n = Anzahl an Knoten \\
		m = Anzahl an Kanten
		\begin{itemize}
			\item Wähle zwei zufällige Knoten aus
			\item Füge Kante zwischen Knoten hinzu, falls neue Kante
			\item Iteriere, bis Anzahl m erreicht
		\end{itemize}
	\end{columns}
	\vspace*{1em}
	Insgesamt sind beide Varianten noch ungerichtet und prüfen nicht auf Kreisfreiheit!
\end{frame}
\begin{frame} {\secname : \subsecname}
	Daher:
	\begin{itemize}
		\item Füge gerichtete, statt ungerichtete Kanten ein
		\item Prüfe beim Hinzufügen ob eine der beiden Kanten $(i,j)$ oder $(j,i)$ bereits existiert
		\item Prüfe, ob es bereits Pfad $j $ nach $i$ gibt mittels rekursiven Tiefensuche
		\item Falls ja, würde neue Kante $e=(i,j) $ Kreis schließen 
		\begin{itemize}
			\item Darf nicht eingefügt werden 
		\end{itemize}
	\end{itemize}
	
\end{frame}	
\subsection{Uniprocessor distribution}
\begin{frame} [fragile]{\secname : \subsecname}
	\begin{lstlisting} {language=Rust, caption=dsfdfs}
fn uunifast(&self, u: f64) -> Vec<f64> {
	assert!(u <= 1.0 && u >= 0.0, "The processor utilization should be percentage between 0 and 1");
	let mut sum_u = u;
	let mut rng = rand::thread_rng();
	let mut vect_u: Vec<f64> = Vec::new();
	for i in 0..self.nodes.len() - 1 {
		let next_sum_u = sum_u * rng.r#gen::<f64>().powf(1.0 / (self.nodes.len() - i) as f64);
		vect_u.push(sum_u - next_sum_u);
		sum_u = next_sum_u;
	}
	vect_u.push(sum_u);
	println!("{:?}", vect_u);
	assert!(u - vect_u.iter().sum::<f64>().abs() < 1e-6);
	vect_u
}
	\end{lstlisting}
	\vspace*{-0.3em}
	\hfill % schiebt den Text nach rechts
	\parbox{0.5\linewidth}{ % Breite der Box (hier halbe Textbreite)
		\raggedleft % Text rechtsbündig innerhalb der Box
		\scriptsize % kleine Schrift
		In uunifast.rs, ab Zeile 34
	}
\end{frame}
\subsection{Uunifast Algorithmus}
\begin{frame} {\secname : \subsecname}
	Constraints für generierte Graphen  \\
	Gilt $\forall i \in \{0..n-1\}$, n = Anzahl der Tasks, i = Index der Task im Schedule
	\begin{itemize}
		\item $c_i  \in \{10..100\} $ \\
		Die Computation time muss für alle Knoten zwischen 10 und 100 liegen.
		\item $d_i \ge c_i + s_i $ \\
		Die Deadline muss spät genug sein, sodass die Task vorher ausgeführt werden kann.
		\item $s_i \ge d_{i-1}, $ wobei $s_0 = 0$ \\
		Eine Task, welche früher gescheduled ist, muss beendet worden sein bevor die nächste startet. Zur Vereinfachung wird $s_i = d_{i-1}$ gesetzt.
		\item $u_i =  \frac{c_i}{t_i}$ \\
		Das Verhältnis aus Computation time und Periode muss der berechneten Utilization entsprechen. Dabei wird der Einfachheit halber $t_i = d_i$ gleichgesetzt.
	\end{itemize}
\end{frame}
\subsection{Dirichlet Rescale Algorithmus}
\begin{frame} {\secname : \subsecname}
\end{frame}
\section{Ergebnisse und Beispiele}
\begin{frame} {\secname }
	\includegraphics [width=0.2 \textwidth] {images/graph.png}
\end{frame}

\end{document}

